\documentclass{blux2017}
\usepackage{mathptmx}
\usepackage[scaled=0.92]{helvet}
\usepackage{courier}

\usepackage{custom}
\usepackage{jsr}

\title{Certifying unstability of Switched Systems using Sum of Squares Programming}
\author{Beno\^it Legat, Rapha\"el M. Jungers\\
        Universit\'e catholique de Louvain\\
        Avenue Georges Lemaitre, 4\\
        B-1348 Louvain-la-Neuve\\
        Belgium\\
    Email: {\tt benoit.legat@uclouvain.be},\\
           {\tt raphael.jungers@uclouvain.be}
        \and
    Pablo A. Parrilo\\
        Laboratory for Information \& Decision Systems\\
        Massachusetts Institute of Technology\\
        77 Massachusetts Avenue\\
        Cambridge MA 02139\\
        USA\\
    Email: {\tt parrilo@mit.edu} }
\date{}

\begin{document}
\maketitle

\section{Introduction}
../../latex/absintro.tex


\section{Joint Spectral Radius and Lyapunov functions}
\label{sec:jsrlyap}

../../latex/jsrdef.tex

If there exists a homogeneous (strictly) \emph{positive} function $f$,
often referred to as \emph{Lyapunov function}, such that
\begin{equation}
  \label{eq:primalconstr}
  f(A_\sigma x) \leq \gamub f(x) \quad \text{ for } \sigma = 1, \ldots, m
\end{equation}
then $\jsr \leq \gamub$.

If we restrict $f$ to be an homogeneous polynomial of degree $2d$ and restrict the postivity conditions on $f$ and on the constraint~\eqref{eq:primalconstr}
to Sum of Squares condition then the search for such function $f$ can be formulated as
a Sum of Squares Program and can be solved using Semidefinite Programming.

Let $\ujsrsos$ be the minimal value of $\gamub$ such that the Sum of Squares Program is feasible.
We have
\begin{equation}
  \label{eq:ujsrsosguarantee}
  \min\Big\{{n+d-1 \choose d}, m\Big\}^{-\frac{1}{2d}} \ujsrsos%
  \leq \jsr \leq \ujsrsos.
\end{equation}

\section{Contributions}
We show that if there exists measures $\mu_\sigma$ for $\sigma = 1, \ldots, m$ such that%
\footnote{$\pushf{A}{\mu}$ is the \emph{pushforward measure} defined to be the measure such that
$\la f, \pushf{A}{\mu} \ra = \la f \circ A, \mu \ra$ for any function $f$.}
\begin{equation}
  \label{eq:dualconstr}
  \sum_{\sigma=1}^m \pushf{A_\sigma}{\mu_\sigma} \geq \gamlb \sum_{\sigma=1}^m \mu_\sigma
\end{equation}
then $\jsr \geq \gamlb$.

If we represent $\mu_\sigma$ by its moments of degree $2d$ and we use the moment relaxation on the existence of a measure with these moments and on the constraint~\eqref{eq:dualconstr} then
the search for such measures $\mu_\sigma$ is the dual of the Sum of Squares Program of \secref{jsrlyap}.

Given measures feasible with $\gamlb$, we provide an algorithm that extracts a trajectory $\sigma_1, \sigma_2, \ldots$ of asymptotic growth rate
\begin{equation}
  \label{eq:useqguarantee}
  \lim_{k \to \infty} \|A_{\sigma_k} \cdots A_{\sigma_2} A_{\sigma_1}\|^{1/k} \geq m^{-\frac{1}{2d}} \gamlb.
\end{equation}
Moreover, if the measures are atomic then with another algorithm, we can extract a periodic trajectory with asymtotic growth rate $\gamlb$.

We can generalize these result when the allowed switching sequences of the system~\eqref{eq:switchsys} are constrained by an automaton.
The guarantees \eqref{eq:ujsrsosguarantee} and \eqref{eq:useqguarantee} holds with $m$ replaced by $\rho(A)$, the spectral radius of the adjacency matrix of the automaton~\cite{philippe2016stability, legat2016generating}.

\bibliographystyle{abbrv}
\bibliography{biblio}

\end{document}
